\section{The companion character is the divided Kummer character}
\label{sec:companion-kummer-equality}

Use the integral compact-homology representation of the degree-\(p\)
cyclic cover and its invariant character \(\widetilde\psi_p\) obtained
from \(b=\epsilon_{-1}\). The possible difference between this character
and the divided Kummer character can be eliminated on full GT.

\begin{theorem}\label{thm:companion-kummer-equality}
For every odd prime \(p\), on the entire multiplier-one subgroup
\(M\) of full profinite GT one has
\begin{equation}
             \boxed{\quad
             \widetilde\psi_p(g)=\frac{\rho_p^I(g)}2\pmod2.
             \quad}                                      
\label{eq:companion-kummer-24}
\end{equation}
More generally, the equality holds for standard positive peripheral
automorphisms \(x\mapsto x\), \(y\mapsto f^{-1}yf\), with
\(f\in F'\), satisfying the two-cycle relation and the integral
Anderson duplication identity. No duplication identity is assumed
for a larger gentle group unless it has been separately established.
\end{theorem}

Use the commutator convention \([x,y]=xyx^{-1}y^{-1}\), and write
\begin{equation}
 [f]=R_g(X,Y)[x,y]\quad\hbox{in }F'/F'',\qquad
 B_g(X,Y)=1+(X-1)(Y-1)R_g(X,Y).
\label{eq:companion-kummer-25}
\end{equation}
Write \(\psi_p=\rho_p^I/2\pmod2\) on \(M\), extending the
earlier notation on \(N\). The coefficient variables record conjugation by \(x,y\). The module
coefficients below are the coefficient-two components. We write
\(\rho_{n,2}^I\) when the integral cocycle is viewed in \(\mathbf Z_2\);
exponents of roots of unity retain the full profinite cocycle.
Relation \eqref{eq:companion-kummer-25} agrees with the normalized Alexander coefficient on \(M\).

\begin{lemma}\label{lem:companion-first-layer}
Let \(O\) be unramified of degree \(2s\) over \(\mathbf Z_2\),
with Frobenius \(\sigma\), involution \(\bar z=\sigma^s z\),
and norm-one group \(T_s\). The companion character is
\(\beta_s(z)=b(v\bar v)\) when \(z=\sigma(v)/v\);
Hilbert 90 permits a unit \(v\), and its relative norm is rational.
Another such unit differs by a rational unit, changing this norm
by a square, so the value is well-defined.
On the first principal-unit layer it is
\begin{equation}
 \beta_s(z)=\operatorname{Tr}_{\mathbf F_{2^s}/\mathbf F_2}
                   \left(\frac{z-1}{2}\bmod2\right)
       \qquad(z\in T_s\cap(1+2O)).                       
\label{eq:companion-kummer-26}
\end{equation}
For an unramified extension of odd degree \(d\),
\begin{equation}
 b(N_{O/\mathbf Z_2}z)
       =\operatorname{Tr}_{\mathbf F_{2^d}/\mathbf F_2}
                   \left(\frac{z-1}{2}\bmod2\right)
       \qquad(z\in1+2O).                                 
\label{eq:companion-kummer-27}
\end{equation}
\end{lemma}

\begin{proof}
For \eqref{eq:companion-kummer-26}, Hilbert 90 gives \(z=\sigma(v)/v\) with \(v\in O^\times\).
Since \(z\equiv1\pmod2\), the residue of \(v\) is fixed by Frobenius
and therefore equals one. Write \(v=1+2x\), and put
\(t=v\bar v\in\mathbf Z_2^\times\). Then
\[
 b(t)=x+\bar x\pmod2,
 \qquad \frac{z-1}{2}=x^2+x\pmod2.
\]
The norm-one condition places the latter value in
\(\mathbf F_{2^s}\), and its trace telescopes to
\(x^{2^s}+x=\bar x+x\). This is \eqref{eq:companion-kummer-26}, by the definition
of \(\beta_s\) through the coinvariant isomorphism.
For \eqref{eq:companion-kummer-27}, expansion of the product of conjugates gives
\(N(1+2x)\equiv1+2\operatorname{Tr}(x)\pmod4\).
\end{proof}

\begin{lemma}\label{lem:cyclic-kummer-fourier}
Let \(g\in M\) be a standard positive peripheral automorphism as
in Theorem~\ref{thm:companion-kummer-equality}. In
\(\mathbf Z_2[C_p]=\mathbf Z_2[X]/(X^p-1)\), put
\[
        C_{p,g}(X)=(X-1)R_g(X,1)=\sum_{j=0}^{p-1}c_jX^j.
\]
Then
\begin{equation}
 C_{p,g}(1)=0,\qquad C_{p,g}(X^{-1})=C_{p,g}(X),
 \qquad \rho_{p,2}^I(g)=-c_0.                            
\label{eq:companion-kummer-28}
\end{equation}
If \(\mathcal S\) contains one representative from each reciprocal
pair of nontrivial \(p\)-th roots of unity, then the sum
\[
                S_{p,g}=\sum_{\zeta\in\mathcal S}C_{p,g}(\zeta)
\]
belongs to \(\mathbf Z_2\), independently of those representatives,
and
\begin{equation}
                 \frac{\rho_{p,2}^I(g)}2=-\frac{S_{p,g}}p.
\label{eq:companion-kummer-29}
\end{equation}
\end{lemma}

\begin{proof}
Use the cyclic \(x\)-cover
\(H_p=\ker(F\to C_p)\), with free basis
\(X_0=x^p\), \(Y_j=x^jyx^{-j}\). The class of \([x,y]\)
in its abelianization is \([Y_1]-[Y_0]\). Conjugation by \(x\)
shifts the indices, and conjugation by \(y\) is trivial.
Consequently \eqref{eq:companion-kummer-25} gives
\[
                    [f]=\sum_j c_j[Y_j]
                  \quad\hbox{in }H_p^{\mathrm{ab}}\otimes\mathbf Z_2.
\]
The zero \(x\)- and \(y\)-exponents of \(f\) give no
\(X_0\)-coordinate and \(\sum_jc_j=0\). By definition
\(\rho_{p,2}^I=-\ell_p(f)_2=-c_0\).

Here is the exact boundary symmetry needed for the second identity
in \eqref{eq:companion-kummer-28}. Write \(f_i=x^ifx^{-i}\). Then
\(\phi(Y_i)=f_i^{-1}Y_if_i\), and
\([f_i]=\sum_jc_{j-i}[Y_j]\). The word
\[
                      Y_0Y_1\cdots Y_{p-1}X_0=(yx)^p
\]
is preserved up to conjugation. Its conjugator can be moved into
\(H_p\) by multiplying by a power of \(yx\). In the free class-two
quotient of \(H_p\), comparison of the \([X_0,Y_i]\) coefficients
makes all coordinates of the boundary conjugator equal. The
remaining coefficient at \([Y_i,Y_j]\) is therefore
\(c_{j-i}-c_{i-j}=0\). This proves \(c_j=c_{-j}\) integrally.
It is the cyclic-cover symmetry, and uses no distribution relation.

Fourier inversion over the unramified splitting ring gives
\[
       p c_0=\sum_{\zeta^p=1}C_{p,g}(\zeta)=2S_{p,g}.
\]
The first equality in \eqref{eq:companion-kummer-28} removes the trivial root, and the
reciprocal symmetry gives the second equality. Frobenius permutes
the reciprocal pairs, so \(S_{p,g}\) is an integral rational
two-adic number. Formula \eqref{eq:companion-kummer-29} follows. In particular \(\rho_p^I\)
is even here without an additional parity assumption.
\end{proof}

\begin{proof}[Proof of Theorem~\ref{thm:companion-kummer-equality}]
The two-cycle relation gives \(R_g(X,Y)=R_g(Y,X)\), and hence
\(B_g(X,Y)=B_g(Y,X)\). The exact duplication identity, in the
completed-group-algebra coordinates used here, is
\begin{equation}
                  B_g(X,X)=X^{2\rho_2^I(g)}B_g(-1,X).
\label{eq:companion-kummer-30}
\end{equation}
This is the integral translation displayed on printed p.~102 in
\cite{Furusho2003}: its exponent is \(-2\Psi_2^{(0)}\), and
\(\rho_2^I=-\Psi_2^{(0)}\). That paper's Theorem~3.1 proves
duplication from (I), (II), (IV), and \((K_2)\). The established
full-GT versions of (IV) and \((K_2)\) therefore supply \eqref{eq:companion-kummer-30}.
We use this explicit integral identity, rather than replacing
the coefficient-two value at \(-1\) by its residue at one.

For a nontrivial \(p\)-th root \(\zeta\), the factor
\(\zeta^{2\rho_2^I(g)}\) has odd order, so every torus character
to \(\mathbf F_2\) kills it. We may compute the companion
character on
\begin{equation}
 z_\zeta=B_g(\zeta,-1)
           =1-2(\zeta-1)R_g(\zeta,-1)
           \equiv1\pmod2.                               
\label{eq:companion-kummer-31}
\end{equation}
This value lies in the required torus: by \eqref{eq:companion-kummer-30} it differs by an
odd-order norm-one root from the compact-homology eigenvalue
\(B_g(\zeta,\zeta)\). Reduction modulo two identifies the
evaluation maps at \(Y=-1\) and \(Y=1\). Thus
\begin{equation}
          \frac{z_\zeta-1}{2}
                   =-C_{p,g}(\zeta)\pmod2.               
\label{eq:companion-kummer-32}
\end{equation}

If \(d=\operatorname{ord}_p(2)=2s\), choose one \(\zeta\)
in each Frobenius orbit. Equation \eqref{eq:companion-kummer-28} places
\(C_{p,g}(\zeta)\) in the degree-\(s\) fixed subfield, and
\[
 S_{p,g}=\sum_{\text{Frobenius orbits}}
       \operatorname{Tr}_{O_0/\mathbf Z_2}C_{p,g}(\zeta).
\]
Equations \eqref{eq:companion-kummer-26} and \eqref{eq:companion-kummer-32} show that
\(\widetilde\psi_p(g)=-S_{p,g}\pmod2\).
If \(d\) is odd, choose one orbit from each reciprocal pair.
Now
\[
 S_{p,g}=\sum_{\text{orbit pairs}}
       \operatorname{Tr}_{O/\mathbf Z_2}C_{p,g}(\zeta),
\]
and \eqref{eq:companion-kummer-27} gives the same conclusion. Finally \(p\) is odd,
so \eqref{eq:companion-kummer-29} gives
\[
          -S_{p,g}\equiv-\frac{S_{p,g}}p
                    =\frac{\rho_p^I(g)}2\pmod2.
\]
This proves \eqref{eq:companion-kummer-24} on all of \(M\).
\end{proof}

\begin{corollary}\label{cor:cubic-two-lifts-equal}
Let \(\tau\) be the earlier coefficient-two cubic frame character.
Then on all of \(M\),
\begin{equation}
                         \xi_{2,3}=\tau+\psi_3.
\label{eq:companion-kummer-33}
\end{equation}
Thus the difference between these two lifts of \(\kappa_{2,3}\)
vanishes on full GT.
\end{corollary}

\begin{proof}
First take \(g\in N\). Proposition~\ref{prop:lr-integral-cubic-beta},
whose scope is \(N\), identifies its coefficient-two
value with
\(T=\varepsilon B_g(\omega,\omega)^{\pm1}\), where
\(\varepsilon\in\{1,-1\}\), \(O=\mathbf Z_2[\omega]\), and
\(T\in1+4O\) has norm one. Both \(\alpha_1\) and \(\beta_1\)
take value one at \(-1\), while
\[
 \alpha_1(T)=\frac{\log T}{4(1+2\omega)}\pmod2=\tau(g),
                     \qquad \beta_1(T)=0.
\]
The possible conjugation or inversion does not change a quadratic
character. Hence \(\xi_{2,3}+\widetilde\psi_3=\tau\) on \(N\).
Theorem~\ref{thm:companion-kummer-equality} proves \eqref{eq:companion-kummer-33} on this subgroup.

We establish the remaining arithmetic values separately. The
compact degree-three cover is
\[
                    C:\quad w^3=t(t-1).
\]
The coordinates \(X=4w\), \(Y=4(2t-1)\) identify its
compactification with
\[
                    E_{16}:\quad Y^2=X^3+16.
\]
The isomorphism with \(E_3:Y^2=X^3+1\), over
\(\mathbf Q(2^{1/3})\), is
\[
                 (X,Y)\longmapsto
                    (X/2^{4/3},Y/4).
\]
Its Galois twisting automorphisms have order dividing three:
changing \(2^{1/3}\) by \(\omega\) multiplies the first
coordinate by \(\omega^{-1}\) and fixes the second.
Likewise, the difference between the restriction of the based
arithmetic action to the cyclic-cover subgroup and the action
using a chosen base point on \(C\) is a deck transformation.
Indeed the two choices differ by conjugation by an element of
\(F\); conjugation by its \(U_3\)-part is trivial on
\(H_1(U_3)\), leaving only \(F/U_3=C_3\).
Both quadratic torus characters kill these order-three
corrections. Consequently, on
\(H^{\mathrm{cyc}}=H\cap M\),
\[
             \xi_{2,3}(\sigma)
                    =\alpha_1(\rho_{E_3,2}(\sigma)).
\]
Here \(\rho_{E_3,2}\) is the Tate representation;
changes between the two CM embeddings do not affect \(\alpha_1\).

Proposition~\ref{prop:adelic-signed-path-comparison} applies on
\(H^{\mathrm{cyc}}\), and gives
\[
 T_\sigma=\kappa(\sigma)\rho_{E_3,2}(\sigma),\qquad
 \kappa(\sigma)=\frac{\sigma(3^{1/4})}{3^{1/4}}
                         =(-1)^{\psi_3(\sigma)}.
\]
Since \(T_\sigma\in1+4O\), applying \(\alpha_1\) gives
\(\tau(\sigma)=\psi_3(\sigma)+\xi_{2,3}(\sigma)\).
Thus \eqref{eq:companion-kummer-33} holds on all of \(H^{\mathrm{cyc}}\), including
its possibly nonzero mass coordinate \(\rho_2^I/2\bmod2\).
Finally \(G=NH\) and \(\lambda(N)=1\) imply
\(M=N H^{\mathrm{cyc}}\). All three terms in \eqref{eq:companion-kummer-33} are
homomorphisms on \(M\), so their equality on these two
factors proves it on \(M\). No full-\(M\) geometric
cubic-frame comparison is used.
\end{proof}

\subsection{The harmonic candidate as an exact comparison of two lifts}

The previous theorem eliminates the companion differences. It does
not identify the coefficient-two character with every other geometric
lift of the same arithmetic target. The harmonic residue supplies
a precise remaining comparison.

For an odd prime \(p\), choose \(i^2=-1\) in
\(\overline{\mathbf F}_p\), put
\(h_p=(p-\chi_{-4}(p))/2\), and define on \(N\)
\[
 (-1)^{\vartheta_{p,4}(g)}=B_{g,p}(i,-1)^{h_p}.
\]
This is the established harmonic residue character. It is the
quadratic torus character of compact homology modulo \(p\) of
the degree-four cyclic cover with monodromies \(1,2,1\).

\begin{proposition}\label{prop:harmonic-two-geometric-lifts}
The transgression of \(\vartheta_{p,4}\) on \(G/N=A\) is
\begin{equation}
                  e_{b,c_p}=b c_p+\frac{p-1}{2}\delta_p.
\label{eq:companion-kummer-34}
\end{equation}
Thus the existing harmonic relative character is exactly
\begin{equation}
                  \Xi_{p,4}
                     =\vartheta_{p,4}-\widetilde\psi_p
                         \quad\hbox{on }N.               
\label{eq:companion-kummer-35}
\end{equation}
Its vanishing on full GT is equivalent to equality of these two
geometric lifts of \eqref{eq:companion-kummer-34}. No vanishing or nonvanishing assertion
for \eqref{eq:companion-kummer-35} is made here.
\end{proposition}

\begin{proof}
The degree-four cover has four punctures and compact genus one.
Its compact deck eigenvalues are \(i,-i\); the order-two
character is peripheral. The positive peripheral action gives
a symplectic similitude normalizer with quotient
\[
                         C_2\times\mathbf F_p^\times.
\]
The two quotient coordinates are the deck exponent modulo four
and the similitude multiplier modulo \(p\), giving characters
\(b,c_p\). The kernel is the split torus \(\mathbf F_p^\times\)
when \(p\equiv1\pmod4\), and the norm-one torus of
\(\mathbf F_{p^2}\) when \(p\equiv3\pmod4\). In both cases
use its unique nonzero quadratic character \(\chi\).

Choose a symplectic lift \(J\) exchanging the two deck eigenlines.
In the split case one may use the matrix with rows \((0,1)\),
\((-1,0)\); in the inert case use \(z\mapsto c\bar z\)
with \(c\bar c=-1\), after normalizing the alternating trace form.
Then \(J^2=-I\). The torus has order \(p-\chi_{-4}(p)\),
which is divisible by four, so \(\chi(J^2)=0\).
Thus the \(b^2\) coefficient vanishes.

In the split case take \(L=\operatorname{diag}(1,t)\) for a
generator \(t\) of \(\mathbf F_p^\times\). Its commutator
with \(J\) has torus parameter \(t\) or \(t^{-1}\), hence
quadratic value one. Also \(L^{p-1}=1\).
In the inert case choose a generator \(v\) of
\(\mathbf F_{p^2}^\times\), and let \(L\) be multiplication
by \(v\). Its norm generates \(\mathbf F_p^\times\).
The commutator has parameter \(v^{p-1}\), or its inverse;
both it and \(L^{p-1}\) have quadratic sign
\(v^{(p^2-1)/2}=-1\) in the norm-one torus.
Hence the mixed coefficient is always one,
and the \(\delta_p\) coefficient is zero in the split case
and one in the inert case. These are all three degree-two
coefficients, proving \eqref{eq:companion-kummer-34} by the homology representation
and naturality of transgression.

The arithmetic transgression is injective, so \eqref{eq:companion-kummer-34} also proves
\(\vartheta_{p,4}|_{K}=\psi_p|_{K}\) directly, without
using an odd Anderson multiplication identity. The difference
therefore extends to a character of \(G\) killing \(H\).
Equation \eqref{eq:companion-kummer-35} follows from \eqref{eq:companion-kummer-24}. Equality of transgression
classes identifies their restrictions arithmetically; it does
not identify the two lifts on \(G\), which is precisely the
remaining relative-kernel issue.
\end{proof}

Equivalently, the unresolved finite compatibility is
\begin{equation}
 B_{g,p}(i,-1)^{(p-\chi_{-4}(p))/2}
       =(-1)^{\chi_b(\rho_{p,2}(g))}
                   \qquad(g\in N).                       
\label{eq:companion-kummer-36}
\end{equation}
The left side uses coefficient \(p\) and the degree-four cover;
the right side uses coefficient two and the degree-\(p\) cover.
The present proof establishes the right side's divided-Kummer
formula exactly. It supplies no additional implication forcing
\eqref{eq:companion-kummer-36}. In particular the exact nonabelian cyclic-filling theorem
cannot be replaced by an endomorphism of the original Alexander
module in an attempt to infer this comparison.
